\begin{textsample}{POS Dim 3 – sp – Score 9.00 – sp/sp\_inf\_curriculo\_e\_ppp.txt}  \label{ex:f3_pos_200}
Concepção de Currículo para Educação \textbf{Infantil} O Parecer CNE/CEB nº 20/2009 afirma que o currículo da Educação \textbf{Infantil} é concebido como um conjunto de práticas que buscam articular as experiências e os saberes das \textbf{crianças} com os conhecimentos que fazem parte do patrimônio cultural, artístico, científico e tecnológico.
Tais práticas são efetivadas por meio de relações sociais que as \textbf{crianças} desde bem \textbf{pequenas} estabelecem com os professores e com as outras \textbf{crianças}, afetando a construção de suas identidades.
No planejamento do currículo devem ser levadas em conta as possibilidades de descobertas, as potencialidades e as genialidades das \textbf{crianças}, mediante o acolhimento genuíno de suas especificidades e interesses singulares.
Isso demanda da instituição de Educação \textbf{Infantil} a promoção de experiências \textbf{lúdicas} e significativas, que de fato permitam às \textbf{crianças} compreenderem e afetarem o mundo no qual estão inseridas.
Assim, faz -se necessário garantir condições para que a \textbf{criança} usufrua do direito de aprender e se desenvolva convivendo, \textbf{brincando}, participando, explorando, expressando e conhecendo -se em contextos culturalmente significativos para ela.
Com isso, a \textbf{creche} e a \textbf{pré-escola} precisam se organizar como espaços de acolhimento, descobertas, interações e \textbf{brincadeira}, com condições que favoreçam o desenvolvimento pleno, num ambiente educativo de qualidade, que contribua significativamente para a construção da aprendizagem de todas as \textbf{crianças}.
Projeto Político Pedagógico Os projetos políticos pedagógicos revelam as concepções e as práticas de cada rede e, mais especificamente, explicitam a identidade da unidade de educação \textbf{infantil} que, presente em um determinado contexto social, deve atender aos anseios da comunidade onde está inserida.
Assim, como ponto de partida, a instituição de Educação \textbf{Infantil} deve construir seu projeto político pedagógico considerando os processos democráticos e participativos, tendo como instrumentos possíveis a avaliação institucional com a participação das \textbf{crianças}, da equipe da escola, das famílias.
Deste modo, os anseios da comunidade escolar são \textbf{acolhidos}, organizados e significados por meio do currículo que, contextualizado, deve contemplar os diferentes tempos, espaços e a cultura local, com vistas a aprofundar as experiências que promovam a aprendizagem e o desenvolvimento das \textbf{crianças}.

% matched lemmas: acolher, brincadeira, brincar, creche, criança, infantil, lúdico, pequeno, pré-escola
\end{textsample}
